\subsection{PI次数4における悪い中心値の有限Ext判定}
\label{subsec:order4-finite-ext-test}
中心上有限性は、自己 $\Ext^1$ の奇数性より強い、全単純加群を尽くす判定へつながる。
以下では代数閉体上へ係数を拡大し、固定した位数4包絡の $U_D$ を扱う。

\begin{proposition}[二つのExt次数と四次元以下の単純加群]
\label{prop:order4-six-seven-test}
$U_D$ が H3 切断を与えるための必要十分条件は、全ての単純右加群 $S,T$ に対し
\[
 \Ext^6_{U_D}(S,T)=\Ext^7_{U_D}(S,T)=0
\]
となることである。ここで全単純加群は有限次元で、その次元は $1,2,3,4$ のいずれかである。
従って無限個の Ext 次数や任意に大きな表現次元を調べる必要はない。
この有限判定の全行列を本稿で既に計算した、という主張ではない。
\end{proposition}
\begin{proof}
$U_D$ は affine 中心上有限の PI次数4の代数なので、代数閉体上の全単純加群は有限次元、
次元は4以下である。これは中心の最大イデアルへ降ろした有限次元代数の単純加群と
PI次数の標準的な表現次元上界による。

有限次元加群 $M$ を、基底空間 $M$ に $x_i$ が $tX_i$ と作用する
$k[t]\otimes M$ という graded $R=B/(\mathcal G_D)$ 加群へ持ち上げる。
$B$ は noetherian AS4、$\mathcal G_D$ は正則中心三次元なので、
Cassidy--Conner--Kirkman--Moore \cite[Theorem 4.7(1),(2)]{CCKMcut}より、
その最小graded自由解消は五段階の後に二周期となる。
$t$ を逆にし次数0を取ると内部shiftは自明化されるので、全有限次元 $M,N$ について
\[
 \Ext^{i+2}_{U_D}(M,N)\simeq\Ext^i_{U_D}(M,N)\qquad(i\ge6)
\]
となる。従って表示した二次数の消滅は、全単純加群の組について次数6以上の消滅を意味する。

中心の最大イデアル $\mathfrak m$ で完備化した環を $\widehat U_{\mathfrak m}$ とする。
これは完備局所可換環上有限なので semiperfect であり、有限生成加群は最小射影解消を持つ。
$\mathfrak m$ 上の単純加群間の Ext は、中心完備化の忠実平坦性により元の Ext と一致する。
最小解消を単純加群へ写すと微分が零となるから、Ext の全組の上記消滅は
各単純加群の射影次元が高々5であることを意味する。
有限次元双対と injective 解消の双対化により
\[
 \Ext^i_U(S,T)\simeq\Ext^i_{U^{\mathrm{op}}}(T^*,S^*)
\]
でもあるので、左単純加群にも同じ上界が成り立つ。

完備環の Jacobson radical を $J$ とする。
左加群 $\widehat U/J$ の射影次元は高々5なので、任意の有限生成右加群の最小射影解消を
$\widehat U/J$ と tensor すると、次数6以上の項は Tor の消滅と Nakayama により零となる。
従って完備化後の全有限生成加群の射影次元は高々5である。
五次 syzygy の射影性は有限中心の全最大点の忠実平坦完備化で検出できるので、
元の $U_D$ にも $\gldim U_D\le5$ が従う。
既出の一affine chart判定により有限大域次元は正確に H3 を意味し、その場合の大域次元は2となる。
逆向きは H3 の $\gldim U_D\le2$ から明らかである。
\end{proof}

\begin{corollary}[悪い中心値を与える有限代数データ]
\label{cor:order4-effective-ext-elimination}
この固定包絡の全 bad $D$ は、$1\le n,m\le4$ の16個の表現次元の組と二つの Ext 次数に対する
有限な多項式階数条件の像として表せる。したがって、有限 Gröbner 計算によって
全 bad 値の monic square-free 多項式を求める、停止する算法が存在する。
この多項式の実際の係数はまだ計算していない。
\end{corollary}
\begin{proof}
固定した $n,m$ に対し、三生成元の $n\times n,m\times m$ 行列に包絡関係と中心切断を課し、
同じ $D$ を持つ表現schemeの積を取る。その座標環を $A$ とする。
$U\otimes_{k[D]}A$ は可換有限型環 $Z(U)\otimes A$ 上有限なので noetherian であり、
$A$ 上平坦である。二つの普遍表現は $A$ 上有限自由である。
一方の普遍表現の有限生成自由解消を第8項まで取ることができる。
各 syzygy も $A$ 上平坦なので、この有限部分の完全性は任意の底変更で保たれる。
もう一方の普遍表現へ Hom を取ると、有限自由 $A$ 加群の有限複体を得る。
その次数6,7の非零cohomology条件は有限個の行列階数条件で表せる。

単純表現の条件も有限な開条件である。
Burnside の定理により、生成行列の単語が全 $\operatorname{Mat}_n$ を張ることと同値であり、
長さ $n^2-1$ までの単語を使えば足りる。
従って単純性の開条件と Ext 非零の閉階数条件を合わせ、余分な変数で開条件の単元を表せば
有限個の affine 代数データとなる。
前命題により、その全像の和集合が正確な bad $D$ 集合である。

この像は Chevalley の定理により $D$ の直線の constructible 部分集合である。
もし無限なら非空開集合を含み、定義体上超越的な $D$ も bad になる。
これは既証の超越的中心値の admissibility に反する。
従って像は有限集合であり、消去による閉包を取っても余分な中心値を加えない。
ここで固定族の有限例外性は、有限中心と有限 Ext 判定からも独立に導かれ、
未証明の全相対 H3 開性を入力にしていない。
最後に radical と monic 化で所望の多項式を得る。
有限中心上の加群表示を用いれば、ここで必要な有限自由解消とsyzygyは
可換有限型環上の有限 Groebner 計算で構成できる。
これは実際の大行列を省略してその計算結果を主張することとは異なる。
\end{proof}
